% leanify-doc.tex -- documentation of the LaTeX package leanify.
% Copyright (C) 2026 Marcus Rockel. Apache License, Version 2.0.
\documentclass[11pt]{article}
\usepackage[T1]{fontenc}
\usepackage{lmodern}
\usepackage[a4paper,margin=30mm]{geometry}
\usepackage{amsmath,amssymb,mathtools}
\usepackage{booktabs,tabularx,array}
\usepackage{listings}
\usepackage{xcolor}
\usepackage[colorlinks=true,allcolors=blue!50!black]{hyperref}
\lstset{basicstyle=\ttfamily\small,columns=fullflexible,keepspaces=true,
  frame=single,framerule=0.3pt,rulecolor=\color{black!30},
  xleftmargin=4pt,xrightmargin=4pt,upquote=true}
\newcommand*\pkg[1]{\textsf{#1}}
\newcommand*\opt[1]{\texttt{#1}}
\newcommand*\cs[1]{\texttt{\textbackslash#1}}

\title{The \pkg{leanify} package\\[2pt]\large A LaTeX article checked by Lean}
\author{Marcus Rockel\\\small\url{https://github.com/Corrram/leanify}}
\date{Version 0.7, 2026-11-10}

\begin{document}
\emergencystretch=2em
\maketitle

\begin{abstract}
Leanify compiles an ordinary mathematical article, written in LaTeX, to the
Lean proof assistant. Every definition and claim of the article becomes a Lean
statement checked by the kernel, and the proofs are checked sentence by sentence
against the results they cite. There is no second, hand-written formal version
that could drift away from the paper: one source is typeset as the article you
submit and compiled by Lean. This document describes the LaTeX package
\pkg{leanify}, the typesetting side of Leanify. The Lean side, the grammar of
the articles it reads and a tutorial are documented at
\url{https://corrram.github.io/leanify/}.
\end{abstract}

\section{What the package does}

Leanify reads the article source itself: its sentences, formulas and the macros
of its preamble. The package does not change what the article says, and Leanify
reads no markup from it. The package
\begin{itemize}
\item loads the packages for the notation the Leanify grammar reads:
  \cs{coloneqq} from \pkg{mathtools}, \cs{mathbb} and \cs{le} from \pkg{amssymb};
\item checks the article while it compiles, when it can (Section~\ref{sec:check});
\item with the option \opt{status}, marks every labelled result in the margin with
  its verification status (Section~\ref{sec:status});
\item records the Lean module that certifies the article (\cs{leanifymodule}).
\end{itemize}
Without options and without shell escape the typeset article is unchanged, so
an article that loads \pkg{leanify} compiles anywhere, for instance on arXiv or
a journal's server.

For example, Leanify compiles the sentence
\begin{lstlisting}
Define $f:(0,\infty)\to\mathbb{R}$ by $f(x)\coloneqq x+\frac1x$.
\end{lstlisting}
to the Lean definition \texttt{f := fun x => x + 1 / x}, with the obligation
that it is well defined, and the theorem
\begin{lstlisting}
\begin{theorem}\label{thm:bound}
For every $x\in(0,\infty)$, $f(x)\ge2$.
\end{theorem}
\end{lstlisting}
to the Lean statement $\forall\,(x:\mathbb{R}),\ 0<x\to 2\le f\,x$, derived in Lean from
exactly the sentences of its proof and the results they cite.

\section{Usage}

\begin{lstlisting}
\usepackage{leanify}            % check as you compile
\usepackage[status]{leanify}    % and mark results in the margin
\usepackage[certify]{leanify}   % and build the certificate first
\leanifymodule{Note.Certified}  % optional: the certifying module
\leanifyreport                  % optional: a one-line summary
\leanifyreadings                % optional: the formal readings
\end{lstlisting}

\noindent The options are:

\medskip
\noindent\begin{tabularx}{\linewidth}{@{}>{\ttfamily}l X@{}}
\toprule
\normalfont Option & Effect \\
\midrule
status & Marks every labelled result in the margin with its verification status,
  read from \texttt{\cs{jobname}.leanify.tex}. \\
check=false & No quick check during the compile. \\
certify & Before the marks are read, updates the pin of the article and builds
  its certificate with \texttt{lake build}. The first time, when the article has
  no Lean files yet, writes them. Implies \opt{status}. \\
vocab=FIELDS & The fields of the article's vocabulary, for instance
  \texttt{vocab=unitsquare} (probability on the unit square, with $[0,1]$ as a type
  of its own) or \texttt{vocab=\{numbers,hyperbolic\}}. The core vocabulary (real
  numbers, integers, their signs) is always there. \\
lean=FILE & The compiled Lean file of the article, if it is not
  \texttt{Compiled.lean} next to it or in the folder that \texttt{leanify init}
  made. \\
exe=COMMAND & The checker to run instead of the project's, for instance
  \texttt{exe=leanify} for an executable on the \texttt{PATH}. With it the quick
  check also runs outside a Lean project. \\
notions=NOUNS & The nouns that the article uses without defining them, for
  instance \texttt{notions=\{graph, finite graph, vertex (vertices)\}}
  (Section~\ref{sec:notions}). Nothing is printed for them. \\
quantities=WORDS & Words that name a real number in the names of quantities that
  the article does not define, for instance \texttt{quantities=\{degree\}}. \\
\bottomrule
\end{tabularx}

\medskip
\noindent The commands are:
\begin{description}
\item[\cs{leanifymodule}\texttt{\{}\textit{module}\texttt{\}}] The Lean module whose
  \texttt{\#leanify\_certify} certifies the article. By default it is
  \texttt{NS.Certified}, for the namespace \texttt{NS} that the compiled Lean file
  compiles.
\item[\cs{leanifyreport}] Prints a one-line summary of the last build where it
  appears, for instance: \emph{12 of 12 claims verified (3 derived in the text from
  their citations); 85 of 90 sentences compiled; 7 definitions well defined}.
\item[\cs{leanifyreadings}] Prints the formal reading of every definition and of
  every statement of the article, with its line, the version of Leanify, the
  vocabulary and the conventions of the grammar: what Lean was asked to prove.
  In an appendix, it lets a reader of the PDF compare each statement with its
  reading. It prints from \texttt{\cs{jobname}.leanify.tex}, like the margin
  marks.
\end{description}

\section{Checking as you compile}\label{sec:check}

When the compiler may run programs (shell escape: \texttt{pdflatex -shell-escape},
\texttt{latexmk -shell-escape}, or the setting of the editor) and the article
lies in a Lean project (a \texttt{lakefile.toml} or \texttt{lakefile.lean} in its
folder or up to three folders above), every compile runs one command,
\texttt{leanify tex}, on the article. It writes its findings as records,
\texttt{\cs{jobname}.leanify.out}, and they become LaTeX
warnings at their lines of the article,
\begin{lstlisting}
Package leanify Warning: g is not defined: sentence not checked on input line
(leanify)                70.
\end{lstlisting}
so editors that read the LaTeX log list them with the other warnings. The
check reads the article with the files it inputs, and a finding in such a file
names its file and line. If some proofs have sentences that Leanify does not
read, one warning says how many proofs are read completely: only their results
are derived from the printed proofs. With
\opt{certify}, the errors of the Lean build and the refuted claims are warnings
too. The full output of Leanify is in \texttt{\cs{jobname}.leanify.log}.
\texttt{lake} must be on the \texttt{PATH} of the compiler; the first check
also builds the checker, which takes a few minutes once. Later compiles run the
checker directly, without \texttt{lake}, and take a fraction of a second. The
records of an unchanged article are the same on every run, so the later passes
of \texttt{latexmk} see no change.

Without shell escape nothing runs, and in a Lean project a warning says so.
Outside a Lean project the package runs nothing and says nothing, unless the
option \opt{exe} names a checker.

\section{Notions that the article does not define}\label{sec:notions}

An article on graphs does not define a graph. If the vocabulary has no field for
its subject, the article names the nouns that it uses without a definition:
\begin{lstlisting}
\usepackage[notions={graph, finite graph, vertex (vertices)},
            quantities={degree}]{leanify}
\end{lstlisting}
A noun is then a kind of object of which nothing is known, and an entry that ends
with a noun (\texttt{finite graph}) is an adjective of it. An irregular plural is
given in parentheses. A sentence such as \emph{Let $\chi(G)$ denote the chromatic
number of a graph $G$.} introduces a symbol for a quantity that the article does
not define. What Lean proves about these notions is proved whatever they mean.

A result whose heading cites its source and that the article does not prove, as in
\verb|\begin{lemma}[\cite{diestel}]|, is \emph{quoted from the literature}. The
article assumes it: it is an axiom of the certificate, and every verified claim
records the quoted results on which it rests. The summary of \cs{leanifyreport}
names what is not defined in the article and how many statements are quoted. A
kind of object is assumed to have an element, so a sentence that says only that
such an object exists is not read as a claim.

\section{Verification status in the margin}\label{sec:status}

With \opt{status}, every labelled result whose status is known gets one of these
marks in the margin:

\medskip
\noindent\begin{tabularx}{\linewidth}{@{}l X@{}}
\toprule
Mark & Meaning \\
\midrule
$\checkmark$\,\textsf{checked} & Every sentence compiled, and the claim is verified in Lean. \\
$\checkmark$\,\textsf{derived} & Every sentence compiled, and the proof in the text
  derives the claim from its citations, each step by the standard tactic, under
  hypotheses that are shown to hold together. \\
$\checkmark$\,\textsf{argued} & As \textsf{derived}, with a reservation that the
  report names: a step of the proof needed a longer Lean argument than the standard
  tactic (the printed proof leaves out more than a routine step there), or the
  hypotheses of a claim are not shown to hold together (it may be true for an empty
  reason). \\
$\circ$\,\textsf{quoted} & The result cites its source and has no proof in the
  article. It is assumed, and the certificate names it among its axioms. \\
$\circ$\,\textsf{partly checked} & The result is verified in part. \\
$\circ$\,\textsf{compiled} & The statement compiled, but is not verified yet. \\
$\times$\,\textsf{false} & The claim is refuted in Lean. \\
\bottomrule
\end{tabularx}

\medskip
The marks are read from \texttt{\cs{jobname}.leanify.tex}, which the Lean command
\texttt{\#leanify\_certify} writes next to the article. If the article is newer
than that file, a warning says that the marks may be out of date. Leave
\opt{status} out for the version you submit.

The summary of \cs{leanifyreport} also says how many steps the standard tactic
closes, and of the sets of hypotheses under which claims are made, how many are
shown to hold together.

\section{Requirements}

The package needs LaTeX 2022-06-01 or newer, \pkg{mathtools} and \pkg{amssymb}.
Checking and the status marks need the Lean toolchain and the Leanify Lean
package in the article's Lake project; typesetting does not. Installation and
a step-by-step tutorial are at \url{https://corrram.github.io/leanify/}.

\section{Licence}

Copyright \copyright\ 2026 Marcus Rockel. The package is released under the
Apache License, Version 2.0.

\end{document}
